\section{Constructive monotone reparametrisation}
\label{app:monotone}

The total coefficient grid $\Theta$ is generated from unconstrained raw
variables rather than by rejecting proposals that violate inequalities.  Define
$\softplus(t)=\log(1+e^t)$ and
$\sigmoid(t)=(1+e^{-t})^{-1}$.  With positive scales
$\lambda_x=e^{\log\lambda_x}$ and $\lambda_Z=e^{\log\lambda_Z}$, construct the
first row and column as
\begin{align}
 \Theta_{11}&=c_0,\\
 \Theta_{1h}&=\Theta_{1,h-1}+\lambda_Z\softplus(a_h),
 &&h=2,\ldots,K_Z,\\
 \Theta_{i1}&=\Theta_{i-1,1}-\lambda_x\softplus(b_i),
 &&i=2,\ldots,K_x.
 \label{eq:boundaries}
\end{align}
The first row therefore increases with $Z$, whereas the first column decreases
with $x$.  Interior points are visited in increasing $i+h$.  For each point,
define
\begin{equation}
 L_{ih}=\Theta_{i,h-1},\qquad U_{ih}=\Theta_{i-1,h},
\end{equation}
and set
\begin{equation}
 \Theta_{ih}=L_{ih}+(U_{ih}-L_{ih})\sigmoid(r_{ih}),
 \qquad i,h\ge2.
 \label{eq:interior}
\end{equation}
If all previously constructed points obey the two monotone orders, then
\begin{equation}
 L_{ih}=\Theta_{i,h-1}\le\Theta_{i-1,h-1}\le
 \Theta_{i-1,h}=U_{ih}.
\end{equation}
Thus Eq.~(\ref{eq:interior}) lies inside a non-empty interval and simultaneously
satisfies
$\Theta_{i,h-1}\le\Theta_{ih}\le\Theta_{i-1,h}$.  Induction over $i+h$
therefore proves the coefficient inequalities in Eq.~(\ref{eq:coefficient-order})
for the complete grid.  Finally, $D=\Theta-\Pi$ is computed deterministically.

For completeness, the derivative identity for a degree-$d_x$ B-spline surface
has the form
\begin{equation}
 \frac{\partial g}{\partial x}=
 \sum_{i=1}^{K_x-1}\sum_{h=1}^{K_Z}
 \rho_{x,i}(\Theta_{i+1,h}-\Theta_{ih})
 \widetilde B_i(x)C_h(Z),\qquad \rho_{x,i}>0,
\end{equation}
where $\widetilde B_i$ is a non-negative degree-$(d_x-1)$ basis function.
The corresponding $Z$ derivative contains
$\rho_{Z,h}(\Theta_{i,h+1}-\Theta_{ih})B_i\widetilde C_h$ with
$\rho_{Z,h}>0$.  The coefficient orders therefore establish the derivative
signs everywhere in the spline domain, including the appropriate one-sided
limits at its boundaries.
